% Why two links per cell are enough: each cell owns the edge to its right
% neighbour and the edge to the one below, so every edge of the grid graph is
% stored exactly once and the boundary cells are where the null links live.
\begin{figure}[htbp]
  \centering
  \begin{tikzpicture}[scale=1.0]
    \foreach \x in {0,1,2} \foreach \y in {0,1,2} {
      \node[cell] (c\x\y) at (1.7*\x, -1.7*\y) {\scriptsize$(\x,\y)$};
    }
    % P0: the link to the right neighbour
    \foreach \x/\xn in {0/1, 1/2} \foreach \y in {0,1,2}
      \draw[link] (c\x\y) -- node[above, font=\tiny, doorgreen] {$P_0$} (c\xn\y);
    % P1: the link to the neighbour below
    \foreach \y/\yn in {0/1, 1/2} \foreach \x in {0,1,2}
      \draw[link] (c\x\y) -- node[right, font=\tiny, doorgreen] {$P_1$} (c\x\yn);
    % the boundary cells own nothing in the direction that leaves the area
    \foreach \y in {0,1,2}
      \draw[link, wallred, dash pattern=on 1.5pt off 1.5pt]
        (c2\y) -- ++(0.95,0) node[right, font=\tiny] {\textsc{nil}};
    \foreach \x in {0,1,2}
      \draw[link, wallred, dash pattern=on 1.5pt off 1.5pt]
        (c\x2) -- ++(0,-0.95) node[below, font=\tiny] {\textsc{nil}};
  \end{tikzpicture}
  \caption{The cell record. Every cell stores two links, $P_0$ to the cell on
  its right and $P_1$ to the cell below it; the cells of the last column and
  the last row carry \textsc{nil} in the direction that leaves the area. Each
  edge of $\gridgraph$ is owned by exactly one cell, so the $2WH-W-H$ edges are
  held without a single one stored twice, and the four-neighbour view a search
  wants is read back in $O(1)$ from the left and upper neighbours' own links.}
  \label{fig:cell-record}
\end{figure}
